\chapter{Селектор относительной скорости семейного синглета и триплета}
\label{ch:v8-baryon-c0-family-triplet-singlet-relative-rate-selector-gate}

После минимального endpoint-расширения connector-носитель распадается на
неэквивалентные семейные компоненты
\begin{equation}
 W_{\rm conn}^{\rm ext}=W_{mathbf1}\oplus W_{mathbf3},
 \qquad \dim W_{mathbf1}=1,
 \qquad \dim W_{mathbf3}=3.
 \label{eq:v8-baryon-c0-singlet-triplet-connector-decomposition}
\end{equation}
Точная commutant-система для симметричной covariance имеет ранг восемь на
десяти неизвестных:
\begin{equation}
 \rank\mathcal C_{\rm cov}=8,
 \qquad \dim\ker\mathcal C_{\rm cov}=2,
 \qquad C=\gamma_1P_{mathbf1}+\gamma_3P_{mathbf3}.
 \label{eq:v8-baryon-c0-singlet-triplet-covariance-commutant}
\end{equation}
Следовательно, физическая семейная симметрия оставляет свободное отношение
\begin{equation}
 r=\frac{\gamma_3}{\gamma_1}>0.
 \label{eq:v8-baryon-c0-singlet-triplet-relative-rate}
\end{equation}

\section{Два точных канонических свидетеля}

Даже после нормировки общего веса до единицы алгебра
$\mathbb C\oplus M_3(\mathbb C)$ имеет центральный симплекс следов
\begin{equation}
 \rho_p=pP_{mathbf1}+\frac{1-p}{3}P_{mathbf3},
 \qquad 0<p<1,
 \qquad \Tr\rho_p=1.
 \label{eq:v8-baryon-c0-singlet-triplet-central-trace-simplex}
\end{equation}
Он реализует весь положительный луч отношений
\begin{equation}
 r(p)=\frac{1-p}{3p},
 \qquad r:(0,1)\longrightarrow(0,\infty).
 \label{eq:v8-baryon-c0-central-trace-relative-rate-orbit}
\end{equation}
Равная скорость на каждую из четырёх стрелок даёт
\begin{equation}
 \rho_{\rm arrow}=\frac14I_4,
 \qquad (\gamma_1,\gamma_3)=\left(\frac14,\frac14\right),
 \qquad r_{\rm arrow}=1.
 \label{eq:v8-baryon-c0-equal-per-arrow-rate-witness}
\end{equation}
Равная полная интенсивность двух неприводимых секторов даёт другой
SO(3)-инвариантный trace-one представитель:
\begin{equation}
 \rho_{\rm sector}=\frac12P_{mathbf1}+\frac16P_{mathbf3},
 \qquad (\gamma_1,\gamma_3)=\left(\frac12,\frac16\right),
 \qquad r_{\rm sector}=\frac13.
 \label{eq:v8-baryon-c0-equal-per-sector-rate-witness}
\end{equation}
Оба положительны и точны, но имеют разные purity и entropy:
\begin{equation}
 \Tr\rho_{\rm arrow}^2=\frac14,
 \quad S(\rho_{\rm arrow})=\log4,
 \qquad
 \Tr\rho_{\rm sector}^2=\frac13,
 \quad S(\rho_{\rm sector})=\frac12\log12.
 \label{eq:v8-baryon-c0-two-rate-witness-purity-entropy}
\end{equation}
Максимум энтропии выбрал бы первый вариант, но сам принцип максимизации
энтропии multiplicity-rate не содержится в текущем parent.

\section{KMS и примитивность не снимают свободу}

При общей боровской частоте KMS связывает прямую и обратную скорости внутри
каждого сектора,
\begin{equation}
 \gamma_i^{\leftarrow}=e^{-\beta\omega}\gamma_i^{\rightarrow},
 \qquad i\in\{1,3\},
 \label{eq:v8-baryon-c0-singlet-triplet-kms-pair-ratios}
\end{equation}
но межсекторное отношение остаётся тем же:
\begin{equation}
 \frac{\gamma_3^{\rightarrow}}{\gamma_1^{\rightarrow}}
 =\frac{\gamma_3^{\leftarrow}}{\gamma_1^{\leftarrow}}=r.
 \label{eq:v8-baryon-c0-kms-does-not-select-relative-rate}
\end{equation}
Для положительных коэффициентов ядро дирахлеевской формы зависит только от
support скачков:
\begin{equation}
 \ker(\gamma_1\mathcal E_{mathbf1}+\gamma_3\mathcal E_{mathbf3})
 =\ker\mathcal E_{mathbf1}\cap\ker\mathcal E_{mathbf3},
 \qquad \gamma_1,\gamma_3>0.
 \label{eq:v8-baryon-c0-primitivity-independent-of-positive-rate-ratio}
\end{equation}
Поэтому примитивность также не выбирает внутреннюю точку положительного
конуса.

\section{Почему полный M4-след недоступен}

На target-блоке grading равна
\begin{equation}
 \Gamma_4=\operatorname{diag}(-1,+1,+1,+1),
 \qquad \dim\{A\in M_4:[A,\Gamma_4]=0\}=10,
 \label{eq:v8-baryon-c0-singlet-triplet-even-endpoint-algebra-dimension}
\end{equation}
то есть максимальная чётная endpoint-алгебра есть
$M_1\oplus M_3$, а не $M_4$. Любая связывающая матричная единица меняет
grading:
\begin{equation}
 [\Gamma_4,E_{01}]=-2E_{01}\ne0.
 \label{eq:v8-baryon-c0-m4-cross-block-grading-obstruction}
\end{equation}
Следовательно, уникальный след полного $M_4$ нельзя использовать без новой
нечётной endpoint-структуры или изменения grading.

Итоговый реестр равен
\begin{equation}
 N_{\rm rate\ selector}=0/6:
 \quad\{SO(3),\ \mathbb C\oplus M_3\text{-trace},\ KMS,
 \text{primitivity},\ s_0a_0\text{-bridge},\ S_{\max}\}.
 \label{eq:v8-baryon-c0-singlet-triplet-rate-selector-ledger}
\end{equation}

\begin{theorem}[Двухвесовой барьер представления $\mathbf1\oplus\mathbf3$]
Семейная симметрия, grading, KMS-отношение и примитивность допускают
произвольное положительное отношение $r=\gamma_3/\gamma_1$. Равная скорость
на стрелку и равная полная скорость на неприводимый сектор дают два разных
точных свидетеля $r=1$ и $r=1/3$.
\end{theorem}

\begin{nogocriterion}[Следовая изотропия не следует из семейной симметрии]
Выбор $\rho=I_4/4$ эквивалентен дополнительной изотропии между
неэквивалентными представлениями $\mathbf1$ и $\mathbf3$. Нельзя считать
$r=1$ физически выведенным только из равенства следовых норм четырёх
стрелок.
\end{nogocriterion}

\begin{protocol}\begin{enumerate}
\item connector-представление: \textbf{$\mathbf1\oplus\mathbf3$};
\item размерность invariant covariance: \textbf{$2$};
\item свободное отношение: \textbf{$r>0$};
\item equal-per-arrow свидетель: \textbf{$r=1$};
\item equal-per-sector свидетель: \textbf{$r=1/3$};
\item KMS выбирает $r$: \textbf{нет};
\item примитивность выбирает $r$: \textbf{нет};
\item полный чётный $M_4$: \textbf{запрещён grading};
\item селекторов: \textbf{$0/6$};
\item следующий гейт: \textbf{parent-origin центрального trace-веса}.
\end{enumerate}\end{protocol}

Машинный сертификат сохранён в
\path{s2t_v8_baryon_c0_family_triplet_singlet_relative_rate_selector_gate_results.json}.